\section{MDLM 的局限性：KV 缓存问题}

\subsection{KV 缓存回顾}

在 AR 模型中，\textbf{KV 缓存}（Key-Value Caching）是推理加速的关键技术。

\begin{knowledgebox}{KV 缓存的工作原理}
在 AR Transformer 中：
\begin{enumerate}
    \item 因为使用因果注意力，每个 token 的激活值只依赖左侧上下文
    \item 已生成 token 的激活值\textbf{永远不会改变}，因为未来的 token 不影响它们
    \item 因此可以缓存（freeze）已计算的 Key 和 Value
    \item 生成新 token 时只需对\textbf{单个 token} 做前向传播，而非整个上下文
\end{enumerate}
这使得 AR 模型即使需要上千步解码，每一步的计算成本都很低。
\end{knowledgebox}

\subsection{为什么 MDLM 不支持 KV 缓存}

MDLM 使用双向注意力，这意味着\textbf{每个 token 的激活值依赖所有其他 token}。

考虑以下场景：
\begin{enumerate}
    \item 初始状态：$[\text{M}, \text{M}, \text{M}, \text{M}]$
    \item 第一步去噪后：$[\text{M}, B, \text{M}, D]$（揭示了位置 2 和 4）
    \item 第二步去噪后：$[A, B, \text{M}, D, \text{M}, F]$（揭示了位置 1 和 6）
\end{enumerate}

在第二步中，虽然 $B$ 和 $D$ 没有改变，但由于新出现了 $A$ 和 $F$，在双向注意力下 $B$ 和 $D$ 的激活值会发生变化。因此\textbf{无法冻结任何位置的激活值}。

\begin{warningbox}{MDLM 的推理瓶颈}
尽管 MDLM 可以用远少于序列长度的步数生成文本（例如用 400-500 步生成 1000 词），但每一步都需要对\textbf{整个序列}做完整的前向传播。对于长序列，注意力的 $O(n^2)$ 复杂度使得每步的成本远高于 AR 模型（带 KV 缓存）的单步成本。因此在实际推理时间上，MDLM 可能反而比 AR 模型更慢。
\end{warningbox}

\subsection{本章小结}

KV 缓存的缺失是 MDLM 最大的实际局限性。双向注意力虽然赋予了模型全局感知能力，但也导致了激活值随输入变化而全局更新，阻碍了 KV 缓存的使用。

\newpage
